\chapter{Аналитическая часть}

В данной лабораторной работе рассматриваются два алгоритма умножения матриц: \textbf{стандартный алгоритм} и \textbf{алгоритм Винограда}.

\section{Стандартный алгоритм умножения матриц}

Стандартный алгоритм умножения матриц используется в линейной алгебре и является базовым методом для перемножения двух матриц. Пусть даны две матрицы \( A \) размером \( n \times m \) и \( B \) размером \( m \times k \), тогда результатом умножения будет матрица \( C \) размером \( n \times k \), где каждый элемент \( c_{ij} \) вычисляется по следующей формуле:

\begin{equation}
    \label{eq:std}
    c_{ij} = \sum_{r=1}^{m} a_{ir} \cdot b_{rj}
\end{equation}

Общая сложность данного алгоритма составляет \( O(n \cdot m \cdot k) \), что делает его неэффективным для больших матриц. В типичном случае при \( n = m = k \) его сложность оценивается как \( O(n^3) \), что может быть ограничивающим фактором при обработке больших данных.\cite{std}

\section{Алгоритм Винограда}

Алгоритм Винограда представляет собой улучшенную версию стандартного алгоритма умножения матриц, который позволяет уменьшить количество операций умножения за счёт предварительного вычисления вспомогательных величин. Пусть даны матрицы \( A \) и \( B \), как и ранее. Для матрицы \( A \) предварительно вычисляются суммы для каждой строки, а для матрицы \( B \) — для каждого столбца:

\begin{equation}
    \label{eq:vin1}
    \text{row\_factor}[i] = \sum_{r=1, r\ \text{нечет}}^{m-1} a_{i,r} \cdot a_{i,r+1}
\end{equation}

\begin{equation}
    \label{eq:vin2}
\text{col\_factor}[j] = \sum_{r=1, r\ \text{нечет}}^{m-1} b_{r,j} \cdot b_{r+1,j}
\end{equation}

Затем элементы матрицы \( C \) вычисляются по формуле:

\begin{equation}
    \label{eq:vin3}
    c_{ij} = -\text{row\_factor}[i] - \text{col\_factor}[j] + \sum_{r=1, r\ \text{нечет}}^{m-1} (a_{ir} + b_{r+1,j}) \cdot (a_{ir+1} + b_{r,j})
\end{equation}

Если \( m \) нечётно, добавляется ещё одно произведение для последнего элемента.

Таким образом, алгоритм Винограда снижает количество операций умножения за счёт предварительных вычислений и замены некоторых умножений сложениями. Это делает его более эффективным для определённых типов задач, особенно когда размер матриц велик. Его сложность в среднем остаётся \( O(n \cdot m \cdot k) \), но константы уменьшаются по сравнению со стандартным алгоритмом.\cite{vin}

\section{Вывод}

Алгоритм Винограда предоставляет более оптимизированный способ умножения матриц по сравнению со стандартным методом, но его эффективность зависит от конкретной задачи и структуры данных.

\clearpage